\chapter{Intégrales generalisees}
\section{Integrale au sens de Riemann}
Soit $f:[a,b]\to \mathbf{R}$ bornee avec $-\infty<a<b<+\infty$.
\begin{definition}[Partition]
 On note une partition de $[a,b]$
$$
P=\{ (x_{0},x_{1}),\dots,(x_{n-1},x_{n}) \}
$$
\end{definition}
\begin{definition}[Sommes de Darboux]
\begin{itemize}
\item Somme de Darboux superieure
$$
\overline{S}(f,P)=\sum_{i=1}^{n} \sup_{x \in [x_{i-1},x_{i}]}f(x)\cdot(x_{i}-x_{i-1}1)
$$
\item Somme de Darboux inferieure
$$
\underline{S}(f,P)=\sum_{i=1}^{n} \inf_{x \in [x_{i-1},x_{i}]}f(x)\cdot(x_{i}-x_{i-1}1)
$$
\end{itemize}
On note $\overline{S}(f)=\inf\{ \overline{S}(f,P)\}$ et $\underline{S}(f)=\sup \{ \underline{S}(f,P) \}$.
\end{definition}

Si $\overline{S}(f)=\underline{S}(f)$ alors $f$ est integrable au sens de Riemann.
\begin{notation}
$$
\int_{a}^{b} f(x) \, dx=\overline{S}(f)=\underline{S}(f) 
$$
\end{notation}
\begin{proposition}
Si $f:[a,b]\to \mathbf{R}$ est continue (par morceaux) alors elle est integrable au sens de Riemann. 
\end{proposition}

On peut generaliser cette notion dans un intervalle qui n'est pas ferme/borne quand $f$ est integrable (au sens de Riemann) sur tout sous intervalle ferme et borne.
\begin{example}
\begin{itemize}
\item $\int_{0}^{1} \ln(x) \, dx$
\item $\int_{0}^{1} \frac{1}{x} \, dx$
\item $\int_{0}^{1} \sin\left( \frac{1}{x} \right) \, dx$
\item $\int_{0}^{\infty}\sin x  \, dx$
\item $\int_{0}^{\infty}e^{ -x }  \, dx$
\item $\int_{1}^{\infty} \frac{1}{x} \, dx$
\item $\int_{-\infty}^{\infty} \frac{1}{1+x^{2}} \, dx$
\end{itemize}
\end{example}
\section{Integrale generalisee sur un intervalle demi-ouvert  borne}
\begin{definition}
Pour $a<b$, soit $f:[a,b)\to \mathbf{R}$ integrable sur tout sous-intervalle ferme et borne de $[a,b)$. Soit
$$
F(t)=\int_{a}^{t} f(x) \, dx, \quad t\in [a,b)
$$
Si $\lim_{ t \to b^{-} }F(t)$ existe alors on dit que $\int_{a}^{b} f(x) \, dx$ existe (converge) et on pose 
$$
\int_{a}^{b} f(x) \, dx=\lim_{ t \to b^{-} } F(t)
$$
Si $\lim_{ t \to b^{-} }F(t)$ n'existe pas (finie) on dit que l'integrale generalisee diverge.
\end{definition}
\begin{remark}
Pour un intervalle $(a,b]$ on prend $\lim_{ t \to a^{+} }F(t)$.
\end{remark}
\begin{example}[Important!]
$\int_{0}^{1} \frac{1}{x^{\alpha}}  \, dx$ avec $\alpha>0$.

Si $\alpha\neq 1$ alors $$F(t)=\int_{t}^{1} \frac{1}{x^{\alpha}} \, dx=\left[ \frac{x^{-\alpha+1}}{1-\alpha} \right]^{1}_{t}=\frac{1}{1-\alpha}(1-t^{-\alpha+1})$$
qu'on evalue quand $t\to0$
$$
\lim_{ t \to 0^{+} } F(t)=\lim_{ t \to 0^{+} } \frac{1}{1-\alpha}(1-t^{-\alpha+1})=\begin{cases}
\frac{1}{1-\alpha}\left( 1-\frac{1}{t^{\alpha-1}} \right)\to \infty & \text{si }\alpha>1 \\
\frac{1}{1-\alpha} & \text{si } \alpha \in(0,1)
\end{cases}
$$
\end{example}
\begin{exercise}
Est-ce que $\int_{0}^{1} \ln x \, dx$ converge?
\end{exercise}
\begin{remark}
Si $f:[a,b)\to \mathbf{R}$ est continue et $\lim_{ t \to b^{-} }f(x)$ existe c-a-d $f$ peut etre prolongee par continuite sur $[a,b]$ alors $\int_{a}^{b} f(x) \, dx$ converge et concide a $\int_{a}^{b} \tilde{f}(x) \, dx$ avec $\tilde{f}$ la fonction prolongee par continuite.
\end{remark}

\section{Integrale generalisee sur un intervalle ouvert borne}
\begin{definition}

Soit $f:(a,b)\to \mathbf{R}$ integrable sur tout sous-intervalle ferme et borne de $(a,b)$. Pour $c\in(a,b)$ si $\int_{a}^{c} f(x) \, dx$ et $\int_{c}^{b} f(x) \, dx$ convergent alors on dit que $\int_{a}^{b} f(x) \, dx$ converge et on pose
$$
\int_{a}^{b} f(x) \, dx =\int_{a}^{c} f(x) \, dx +\int_{c}^{b} f(x) \, dx 
$$
sinon $\int_{a}^{b} f(x) \, dx$ diverge.
\end{definition}
\begin{remark}
Cette definition ne depend pas du choix de $c$. Si $\int_{a}^{b} f(x) \, dx$ converge pour $c\in(a,b)$ alors si on prend $\tilde{c}\in(a,b)$ alors
$$
\int_{a}^{\tilde{c}} f(x) \, dx=\int_{a}^{c} f(x) \, dx +\int_{c}^{\tilde{c}} f(x) \, dx
$$
\end{remark}

\begin{lemma}[Critere de comparaison]
Soient $f,g\in[a,b)\to \mathbf{R}$ deux fonctions integrables sur tous sous-intervalles fermes et bornes de $[a,b)$. S'il existe $c\in[a,b)$ tel que $\forall x \in [c,b)$ on a $0\leq f(x)\leq g(x)$ et $\int_{a}^{b} g(x) \, dx$ converge alors $\int_{a}^{b}f(x)  \, dx$ converge aussi.

Si $\int_{a}^{b} f(x) \, dx$ diverge alors $\int_{a}^{b}g(x)  \, dx$ diverge aussi.
\end{lemma}
\begin{proof}
Soit $F(t)=\int_{a}^{t} f(x) \, dx$ croissante et bornee supérieurement alors $\lim_{ t \to b^{-} }F(t)$ existe et est finie.
\end{proof}
\begin{definition}
Soit $I$ un intervalle de la forme $[a,b),(a,b),(a,b]$ et $f:I\to \mathbf{R}$ integrable sur tout sous-intervalles ferme et borne.

On dit que l'integrale generalisee converge absolument si
$$
\int_{a}^{b} |f(x)| \, dx
$$
converge.
\end{definition}
\begin{remark}
Si $\int_{a}^{b} f(x) \, dx$ converge absolument alors elle converge. 
\end{remark}